On considère le système chimique obtenu en mélangeant les quantités de matière
suivantes~: ${n_{1}=\qty{1.5e-2}{\mol}}$ d'acide nitreux $\ce{HNO2}$,
$n_{2}=\qty{3e-2}{\mol}$ de méthanoate de sodium $\ce{Na^+_{(aq)} +
HCOO^-_{(aq)}}$, $n_{3}=\qty{3e-2}{\mol}$ de nitrite de sodium $\ce{Na^+_{(aq)}
+ NO^-_2_{(aq)}}$ et $n_{4}=\qty{1.5e-2}{\mol}$ d'acide méthanoïque
$\ce{HCOOH}$. Soit $V$ le volume total du mélange réactionnel.

L'équation de la réaction entre l'acide nitreux $\ce{HNO2}$ et les ions
méthanoate $\ce{HCOO^-_{(aq)}}$ s'écrit~:
\[
    \ce{HNO2_{(aq)} + HCOO^-_{(aq)} <=>[(1)][(2)] NO^-_2_{(aq)} + HCOOH_{(aq)}}
\]

\begin{donnees}
    $\mathrm{pK_{A_{1}}}=
    \mathrm{pK_{A}}(\ce{HNO2_{(aq)}}\ttslash{}\ce{NO^-_2_{(aq)}})=3,2$~; \quad
    $\mathrm{pK_{A_2}}=
    \mathrm{pK_{A}}(\ce{HCOOH_{(aq)}}\ttslash{}\ce{HCOO^-_{(aq)}})=3,8$
\end{donnees}

\begin{enumerate}
    \item Déterminer la valeur du quotient de réaction $Q_{r,i}$ à l'état
          initial du système chimique.
    \item Montrer que la constante d'équilibre $\mathrm{K}$ associée à
          l'équation chimique précédente s'écrit~:
          ${\mathrm{K}=10^{(\mathrm{pK_{A_2}-pK_{A_1}})}}$.
          Calculer la valeur de $\mathrm{K}$.
    \item Indiquer, en justifiant, dans quel sens évolue spontanément le système
          chimique à partir de son état initial.
\end{enumerate}


